% ===========================================================================
% Adaptive SIMEX for Poisson Regression with Classical Measurement Error:
% Theory, Methods, and Applications  ---  REVISED MANUSCRIPT (JoE format)
% Gaurav Garg, Ramakrushna Mishra
% elsarticle class; compile pdflatex x2
% ===========================================================================
\documentclass[3p,11pt]{elsarticle}
\journal{Journal of Econometrics}
\biboptions{authoryear,round}
\usepackage{amsmath,amssymb,amsthm,mathtools,bm}
\usepackage{booktabs,graphicx}
\usepackage[colorlinks=true,linkcolor=blue,citecolor=blue]{hyperref}
\graphicspath{{../results/}}

\theoremstyle{plain}
\newtheorem{theorem}{Theorem}
\newtheorem{lemma}{Lemma}
\newtheorem{corollary}{Corollary}
\theoremstyle{definition}
\newtheorem{assumption}{Assumption}
\newtheorem{remark}{Remark}

\newcommand{\E}{\mathbb{E}}
\newcommand{\R}{\mathbb{R}}
\newcommand{\N}{\mathbb{N}}
\newcommand{\Prob}{\mathbb{P}}
\newcommand{\Var}{\mathrm{Var}}
\newcommand{\Cov}{\mathrm{Cov}}
\newcommand{\bet}{\bm{\beta}}
\newcommand{\bu}{\bm{u}}
\newcommand{\bw}{\bm{w}}
\newcommand{\bW}{\bm{W}}
\newcommand{\bX}{\bm{X}}
\newcommand{\bY}{\bm{Y}}
\newcommand{\beps}{\bm{\varepsilon}}
\newcommand{\bSig}{\bm{\Sigma}}
\newcommand{\bPhi}{\bm{\Phi}}
\newcommand{\bphi}{\bm{\phi}}
\newcommand{\btH}{\bm{H}}
\newcommand{\btV}{\bm{V}}
\newcommand{\bI}{\bm{I}}
\DeclareMathOperator*{\argmax}{arg\,max}
\DeclareMathOperator*{\argmin}{arg\,min}
\providecommand{\cmark}{\checkmark}
\providecommand{\xmark}{$\times$}
\newcommand{\asimex}{\textup{A-SIMEX}}

\begin{document}

\begin{frontmatter}

\title{Adaptive SIMEX for Poisson Regression with Classical
Measurement Error: Theory, Methods, and Applications}

\author[iiml]{Gaurav Garg\corref{cor1}}
\ead{ggarg@iiml.ac.in}
\author[iiml]{Ramakrushna Mishra}
\affiliation[iiml]{organization={Decision Sciences Area,
  Indian Institute of Management Lucknow},
  addressline={Prabandh Nagar},
  city={Lucknow}, postcode={226013}, state={Uttar Pradesh},
  country={India}}
\cortext[cor1]{Corresponding author.}

\begin{abstract}
Poisson regression is widely used for count outcomes, but covariates are
frequently observed with classical measurement error, and the resulting
naive estimates are attenuated with invalid confidence intervals. We
develop Adaptive SIMEX (A-SIMEX), which automates the two ad-hoc choices
of SIMEX, namely the extrapolation region and the polynomial degree, by
cross-validation with a one-standard-error parsimony rule, adds a
ridge-corrected second extrapolation stage, and is accompanied by
complete asymptotic theory: consistency under a growing-degree regime
that we make explicit; asymptotic normality with an exact decomposition
of the asymptotic variance into data and simulation components (the
latter vanishing at rate $B^{-1}$ in the number of simulation
replicates); and efficiency results relative to the oracle estimator,
including an information-monotonicity lemma showing that no estimator
based on noisy covariates can improve on the oracle to first order,
while the second-order efficiency comparison is signed neither way. The
theory corrects several claims of earlier SIMEX treatments and resolves
the tension between consistency (which requires the polynomial degree to
grow like $\log n$) and inference through an explicit window condition
on the degree growth rate. In simulation studies with $1{,}000$ Monte
Carlo repetitions per design cell, including misspecified (Laplace)
error, negative-binomial outcomes, error variance estimated from
replicated measures, and corrected-score and regression-calibration
competitors, A-SIMEX removes most of the naive attenuation with valid
inference when the error is moderate, and we identify an important
limitation of automatic tuning: held-out prediction on the observed
scale systematically under-extrapolates when the error is large,
because its population minimizer is the naive pseudo-parameter. We
characterize the resulting bias--variance trade-off, quantify the
efficiency premium of extrapolation against the oracle, and give
practical remedies. The methods are implemented in an open-source R
package with an analytic sandwich variance estimator that we validate
against the parametric bootstrap, and are applied to NHANES 2017--2018,
where the measurement error is directly measurable: self-reported
weight in a model for the count of prescription medications. All
results are reproducible from a documented archive with fixed seeds.

\end{abstract}

\begin{keyword}
measurement error \sep SIMEX \sep Poisson regression \sep asymptotic
theory \sep generalized linear models \sep replication archive
\JEL I13 \sep C21 \sep C25
\end{keyword}

\end{frontmatter}

% ---------------------------------------------------------------------------
\section{Introduction}\label{sec:intro}

Economists measure the covariates they regress on, and the
measurements are frequently wrong. In labor and health economics,
self-reported income, earnings, and weight differ from administrative
records in well-documented ways, and the resulting demand and health
production estimates inherit the discrepancy
\citep{griliches1986,bound2001}. In demand systems estimated from
household surveys, reported consumption shares are noisy versions of
the expenditures that economic theory needs, so price and income
elasticities are biased in unknown directions. In policy evaluation,
treatment intensity and exposure are measured by program records or
monitors that misreport systematically, misclassifying the very
variable whose effect is of interest
\citep{mahajan2006,hu2008}. Many of the resulting outcome variables,
from doctor visits to patent counts to crime incidents, are counts,
and the Poisson regression model with its overdispersed extensions is
the standard tool for them \citep{mullahy1986,cameron2013,
winkelmann2008}. A long econometric literature documents that
plugging mismeasured covariates into nonlinear models attenuates
estimates, inflates variances, and destroys interval coverage
\citep{hausman2001,bound2001}; Section~\ref{sec:sims1} quantifies the
damage cell by cell for the Poisson case.

In nonlinear models the problem is structurally different from the
familiar regression-dilution algebra of the linear model. The
pseudo-true parameter of the naive fit is defined only implicitly,
through a nonlinear moment condition that mixes the error variance
with the moment-generating function of $X$ evaluated at the parameter
itself; the attenuation is therefore not a simple constant and is not
sign-monotone in all directions
(Lemma~\ref{lem:curve}(iii) computes it exactly to first order).
Classical measurement error in GLMs is the setting of a substantial
statistics literature \citep{carroll2006}, and econometric work has
developed semiparametric and identification-based alternatives
\citep{lewbel1997,schennach2004,mahajan2006,hu2008}; yet for the
applied count regression the practical menu remains short, and
simulation-based corrections of the SIMEX type are far better known in
biostatistics than in economics \citep{nakamura1990,cook1994}. This
paper is written for the applied econometrician who has a count
outcome, a mismeasured regressor, and either an external estimate of
the error variance or replicated measurements, and who needs a
correction that is automatic, documented, and honest about its own
limits.

Simulation extrapolation (SIMEX; \citealp{cook1994,stefanski1995}) is
arguably the most practical member of the correction toolkit: it makes
no model for the distribution of $X$, requires only the error
variance, reduces the correction to a sequence of ordinary GLM fits,
and extends mechanically to misclassification
\citep{kuechenhoff2006} and semiparametric settings
\citep{apanasovich2009}. Yet SIMEX as usually practiced leaves the two
consequential choices, the extrapolation region and the polynomial
degree, to the analyst, offers no rigorous asymptotic theory for
GLM settings, and provides no closed-form standard errors. Table~\ref{tab:gap}
summarizes where the available approaches stand against the
desiderata that motivate this paper: automatic tuning, formal
asymptotics, usable inference machinery, freedom from validation
data, and scalability to moderate dimensions. The table makes clear
that no existing method occupies the full row that an applied
econometrician needs; it also previews our claim that A-SIMEX fills
the row, with one honest asterisk, namely that its automatic
tuning inherits a documented under-extrapolation bias when the
measurement error is large (Section~\ref{sec:sims1}).

\begin{table}[t]
\centering\footnotesize
\setlength{\tabcolsep}{4pt}
\caption{Available approaches to Poisson (GLM) regression with
classical covariate measurement error, against the desiderata
motivating this paper. ``Automatic'' = no analyst-chosen
smoothing/extrapolation tuning; ``Inference'' = analytic standard
errors beyond point estimation. Superscripts flag caveats discussed in
the text and quantified in Section~\ref{sec:sims}.}
\label{tab:gap}
\begin{tabular}{p{5.6cm}ccccc}
\toprule
 & Automatic & Consistency & Inference & No validation & Moderate-$p$ \\
 &  & theory & machinery & data required & scalable \\
\midrule
Naive MLE \citep{fuller1987} & \cmark & $^{\mathrm{a}}$ & \xmark & \cmark & \cmark \\
Regression calibration \citep{prentice1982,rosner1989} & \cmark & $^{\mathrm{b}}$ & approximate & \cmark & \cmark \\
Corrected score \citep{nakamura1990} & \cmark & \cmark & \cmark & \cmark & $^{\mathrm{c}}$ \\
Deconvolution / functional \citep{fan1991,delaigle2008} & \xmark & \cmark & \xmark & \cmark & \xmark \\
Validation likelihood \citep{qin2000,chen2011} & \cmark & \cmark & \cmark & \xmark & \cmark \\
Bayesian hierarchical ME \citep{richardson1993,berry2002} & \cmark & $^{\mathrm{d}}$ & posterior & \xmark & \xmark \\
Standard SIMEX \citep{cook1994,stefanski1995} & \xmark & \xmark & bootstrap only & \cmark & \cmark \\
\midrule
\asimex{} (this paper) & \cmark$^{\mathrm{e}}$ & \cmark & \cmark & \cmark & \cmark \\
\bottomrule
\end{tabular}

{\footnotesize\emph{Notes.} SIMEX extensions to misclassification and
semiparametric settings: \citet{kuechenhoff2006} and
\citet{apanasovich2009}, respectively. Deconvolution rates:
\citet{fan1993}.
$^{\mathrm{a}}$consistent for the naive pseudo-parameter $\beta(1)$,
not for $\bet_0$ (Lemma~\ref{lem:curve}).
$^{\mathrm{b}}$approximate bias correction whose error is first order
in the error variance (Section~\ref{sec:sims1}).
$^{\mathrm{c}}$RMSE deteriorates sharply when many coefficients are
weakly identified (Section~\ref{sec:sims3}).
$^{\mathrm{d}}$posterior consistency known in special cases.
$^{\mathrm{e}}$the CV criterion systematically under-extrapolates when
the error is large, a limitation we document and quantify rather
than suppress (Section~\ref{sec:sims1}).}
\end{table}

This paper develops Adaptive SIMEX (A-SIMEX) and subjects it to a
level of scrutiny that we argue should be standard for
bias-correction methods. Our contributions are as follows.

\begin{enumerate}\itemsep3pt
\item \textbf{Automated hyperparameter selection, with its limits
      quantified.} Extrapolation region and polynomial degree are
      selected by $k$-fold cross-validation of the held-out Poisson
      deviance, with a one-standard-error parsimony rule. This removes
      subjectivity and is fully reproducible. Our simulations,
      however, reveal a structural limitation that we do not paper
      over: the population minimizer of the observed-scale CV criterion
      is the naive pseudo-parameter $\beta(1)$, not $\beta_0$, so CV
      systematically under-extrapolates when the measurement error is
      large (Section~\ref{sec:sims1}). Automatic tuning is a solution
      for small-to-moderate error and a documented failure mode for
      large error, for which we propose deconvolution-aware remedies.
\item \textbf{A bias-corrected two-stage extrapolation and an exact
      variance estimator.} The primary polynomial fit is stabilized by
      a ridge second stage, and inference uses an analytic sandwich
      that decomposes the asymptotic variance exactly into a
      \emph{data} component and a \emph{simulation} component of order
      $B^{-1}$; this decomposition corrects the additive
      ``parameter-plus-extrapolation'' split that appears in earlier
      SIMEX treatments, and it extends robustly to overdispersed
      outcomes and to estimated error variance
      (Sections~\ref{sec:sims-nb} and~\ref{sec:sims-sigma2}).
\item \textbf{Complete and corrected asymptotic theory.} We prove
      consistency and asymptotic normality of the A-SIMEX estimator, a
      growing-degree central limit theorem that reconciles the
      $\log n$ degree growth needed for consistency with centered
      inference through an explicit window condition on the
      analyticity margin, and oracle-efficiency results comprising an
      information-monotonicity lemma (no estimator based on noisy
      covariates beats the oracle to first order) and an explicit
      second-order expansion whose sign is design-dependent. Along the
      way we correct three claims that circulate informally for SIMEX:
      the exactness of the extrapolation target, the form of the
      variance decomposition, and the sign of the second-order
      efficiency loss.
\item \textbf{Evidence at replication standard.} All simulations
      ($1{,}000$ Monte Carlo repetitions per cell) and the application
      are generated by a single seed-fixed pipeline; the paper reports
      what the pipeline produces, including results that are
      unflattering to the method.
\end{enumerate}

The methods are implemented in an open-source R package, and every
number in the paper is regenerable from a documented archive with
fixed seeds. The plan is as follows. Section~\ref{sec:background}
reviews the model and standard SIMEX. Section~\ref{sec:method}
presents A-SIMEX: cross-validated selection of
$(\lambda_{\max},\text{degree})$ (Section~\ref{sec:cv}), the ridge
second stage (Section~\ref{sec:ridge}), and the analytic sandwich
variance estimator (Section~\ref{sec:var}). Section~\ref{sec:theory}
states the main results: the population SIMEX curve and its
analyticity (Lemma~\ref{lem:curve}), the consistency theorem
(Theorem~\ref{thm:cons}), asymptotic normality with the exact
variance decomposition (Theorem~\ref{thm:an},
Corollary~\ref{cor:var}), the growing-degree central limit theorem
(Theorem~\ref{thm:grow}), and the oracle-efficiency results
(Lemma~\ref{lem:info}, Theorem~\ref{thm:are}). Section~\ref{sec:sims}
reports the simulation studies; Section~\ref{sec:app} the application
to NHANES 2017--2018; Section~\ref{sec:software} the software;
Section~\ref{sec:discussion} limitations and extensions. Complete
proofs are in the Supplementary Material.

\label{sec:contrib}\label{sec:roadmap}

% ---------------------------------------------------------------------------
\section{Background and Setup}\label{sec:background}

\subsection{Poisson Regression}\label{sec:poisson}

We observe $n$ independent subjects $i=1,\dots,n$ with count outcome
$Y_i\in\N_0$ and covariate vector $X_i\in\R^p$. The Poisson regression
model specifies
\begin{equation}\label{eq:model}
  Y_i\mid X_i\;\sim\;\mathrm{Poisson}(\mu_i),
  \qquad \mu_i=\exp\bigl(X_i^\top\bet_0\bigr),
\end{equation}
where $\bet_0\in\Theta\subset\R^p$ is the parameter of interest, with
log-likelihood
\begin{equation}\label{eq:loglik}
  \ell_n(\bet)=\sum_{i=1}^{n}\Bigl\{Y_i\,X_i^\top\bet
  -\exp\bigl(X_i^\top\bet\bigr)\Bigr\}.
\end{equation}
The oracle maximum likelihood estimator
$\hat\beta_{\mathrm{oracle}}$, infeasible because it requires
$X_i$, solves $\nabla\ell_n(\bet)=\bm 0$ and is asymptotically
normal with variance equal to the inverse Fisher information.

\subsection{Classical Measurement Error Model}\label{sec:me}

In practice we observe only a noisy version,
\begin{equation}\label{eq:cme}
  W_i=X_i+\beps_i,
\end{equation}
with $\beps_i\sim N(\bm 0,\bSig_\epsilon)$ independent of $(X_i,Y_i)$
and $\bSig_\epsilon=\operatorname{diag}
(\sigma^2_{\epsilon,1},\dots,\sigma^2_{\epsilon,p})$ assumed known or
consistently estimable from auxiliary data (validation subsamples or
replicated measurements; Section~\ref{sec:sims-sigma2} implements the
latter).

\begin{assumption}[Classical error structure]\label{ass:cme}
$\beps_i\sim N(\bm 0,\bSig_\epsilon)$ independently across $i$,
independent of $(X_i,Y_i)$; $\bSig_\epsilon$ is known (or
consistently estimated).
\end{assumption}

\subsection{The Naive Estimator and Its Bias}\label{sec:naive}

Fitting the Poisson GLM with $W_i$ in place of $X_i$ yields the naive
MLE $\hat\beta_{\mathrm{naive}}$, which is attenuated toward zero. The
attenuation is characterized exactly in
Lemma~\ref{lem:curve}(iii) below: the pseudo-true parameter of the
naive fit is $\beta(1)=\bet_0+\tilde\bet+O(\|\bSig_\epsilon\|^2)$,
with first-order bias
\begin{equation}\label{eq:naivebias}
  \tilde\bet
  =-\bm I_X(\bet_0)^{-1}\Bigl[
    \tfrac12\bigl(\bet_0^\top\bSig_\epsilon\bet_0\bigr)\,
    \E\bigl(e^{X_i^\top\bet_0}X_i\bigr)
    +\bSig_\epsilon\bet_0\,\E\bigl(e^{X_i^\top\bet_0}\bigr)\Bigr],
\end{equation}
where $\bm I_X(\bet_0)=\E\bigl[e^{X_i^\top\bet_0}X_iX_i^\top\bigr]$ is
the oracle information. For $p=1$, $X\sim N(0,1)$ this gives
$\tilde\beta=-\sigma^2_\epsilon\beta_0(1+\beta_0^2/2)/(1+\beta_0^2)<0$
for $\beta_0>0$: attenuation, with the same first-order structure the
linear model is famous for but governed by the moment-generating
function of $X$ at $\bet_0$ rather than by a closed-form dilution
factor.

\subsection{Standard SIMEX}\label{sec:stdsimex}

Simulation extrapolation \citep{cook1994} augments the observed error
artificially and extrapolates backward:

\begin{center}
\small
\begin{tabular}{ll}
\toprule
& \textbf{Algorithm 1 (Standard SIMEX)}\\
\midrule
\multicolumn{2}{l}{\textbf{Require:} data $(Y_i,W_i)$; assumed
$\bSig_\epsilon$; grid $\Lambda=\{\lambda_0=0<\cdots<\lambda_K\}$;
$B$ replicates.}\\
1 & \textbf{for} each $\lambda\in\Lambda$, $b=1,\dots,B$:\\
2 & \quad draw $Z_{ib}\sim N(\bm 0,\bSig_\epsilon)$ i.i.d., set
$W_i^{(b)}(\lambda)=W_i+\sqrt{\lambda}\,Z_{ib}$;\\
3 & \quad fit the Poisson GLM of $Y_i$ on
$W_i^{(b)}(\lambda)$, obtaining $\hat\beta_b(\lambda)$;\\
4 & \quad average: $\bar\beta(\lambda)=B^{-1}\sum_b\hat\beta_b(\lambda)$.\\
5 & \textbf{Extrapolate:} fit $p(\lambda)=a+b\lambda+c\lambda^2$ to
$(\lambda_j,\bar\beta(\lambda_j))$ and report
$\hat\beta_{\mathrm{SIMEX}}=p(-1)$.\\
\bottomrule
\end{tabular}
\end{center}

The intuition: increasing $\lambda$ compounds the error, moving
$\bar\beta(\lambda)$ further from $\bet_0$ along a smooth bias curve;
extrapolating to $\lambda=-1$ ``removes'' the original error.
Section~\ref{sec:population} makes this precise: the population curve
$\beta(\lambda)$ depends on the error only through the \emph{total}
variance $(1+\lambda)\bSig_\epsilon$, is real-analytic in $\lambda$,
and satisfies $\beta(-1)=\bet_0$ exactly, but the naive reading
``$W(-1)=W-\epsilon=X$'' is wrong (the simulated noise is independent
of the observed error and cannot be subtracted), a distinction that
matters for every formal claim.

\subsection{Limitations of Standard SIMEX}\label{sec:limitations}

\begin{enumerate}\itemsep2pt
\item \emph{Ad-hoc hyperparameters:} the grid $\Lambda$ and the
      polynomial degree are chosen by hand, with no principled
      guidance and consequences for reproducibility.
\item \emph{Extrapolation instability:} polynomial extrapolation from
      $[0,\lambda_{\max}]$ to $\lambda=-1$ amplifies noise (by a
      factor governed by the weights $\sum_j|c_j|$; see
      Theorem~\ref{thm:are}), and high-degree polynomials oscillate
      outside the fitted range.
\item \emph{No formal asymptotics:} the bias, variance, and coverage
      properties of $\hat\beta_{\mathrm{SIMEX}}$ had not been
      rigorously characterized for GLMs; existing analyses are
      heuristic.
\item \emph{Computational cost:} $B\times|\Lambda|$ GLM fits per
      evaluation, burdensome when $p$ is moderate, though
      The Supplementary Material (Section~S2) shows that batching and warm
      starts reduce this to seconds.
\end{enumerate}

\subsection{Related Work}\label{sec:related}

Corrected score equations \citep{nakamura1990} replace the score by an
estimate whose conditional expectation given $(X,Y)$ is the true
score; they require the error distribution (as SIMEX does) rather than
a model for $\Pr(X\mid W)$, but their variance explodes when many
coefficients are weakly identified (Section~\ref{sec:sims3}).
Deconvolution and functional methods \citep{fan1993,delaigle2008}
estimate the density of $X$ from $W$ nonparametrically; they are
elegant but ill-posed and hard to scale. Empirical likelihood
\citep{qin2000} can incorporate validation data but requires the
validation sample itself. Our empirical comparisons include the first
two families (Sections~\ref{sec:sims1}--\ref{sec:sims3}); the variance
estimation from replicated measures covers the third's typical data
structure.

% ---------------------------------------------------------------------------
\section{Adaptive SIMEX: Method and Implementation}\label{sec:method}

\subsection{Data-Driven Hyperparameter Selection via
Cross-Validation}\label{sec:cv}

The core innovation of A-SIMEX is to treat the extrapolation region
and polynomial degree as tuning parameters selected by $k$-fold
cross-validation using prediction loss on held-out data.

\subsubsection*{Cross-validation objective}

Partition the data into $K_{\mathrm{fold}}$ folds (typically five). For
a candidate pair $(\lambda_{\max},p_{\mathrm{deg}})$:

\emph{Fit stage.} For fold $k$, using data excluding fold $k$, run
Algorithm~1 with extrapolation grid $\{0,0.5,1,\dots,\lambda_{\max}\}$ and
degree $p_{\mathrm{deg}}$, obtaining the extrapolated estimate
$\hat\beta^{(-k)}$. Within a fold, the curve on the union grid is
fitted once and each candidate $\lambda_{\max}$ is evaluated on the
corresponding prefix (the \emph{prefix property}), with noise
draws shared across $\lambda$ within a replicate.

\emph{Prediction stage.} On the held-out fold, compute the Poisson
deviance of the fitted counts
$\hat Y_{i,k}=\exp(W_{i,k}^\top\hat\beta^{(-k)})$:
\begin{equation}\label{eq:cv}
  D_k(\lambda_{\max},p_{\mathrm{deg}})
  =2\sum_{i\in\text{fold }k}\Bigl[
    Y_i\log\frac{Y_i}{\hat Y_{i,k}}-\bigl(Y_i-\hat Y_{i,k}\bigr)\Bigr],
\end{equation}
with the convention $0\log0=0$, and aggregate
$\mathrm{CV}(\lambda_{\max},p_{\mathrm{deg}})
=K_{\mathrm{fold}}^{-1}\sum_kD_k$. The selected pair minimizes
$\mathrm{CV}$; when candidates lie within one standard error of the
minimum, the \emph{parsimonious} choice (smaller degree, then smaller
$\lambda_{\max}$) is preferred. Section~\ref{sec:sims1} quantifies a
bias of this criterion toward under-extrapolation and
Section~\ref{sec:sims-val} discusses remedies; the selected pair is
additionally reported in all output so that practitioners can inspect
the choice.

\subsubsection*{Computational efficiency}

The CV procedure is parallelizable at two levels (across folds; across
grid points), and the prefix property reduces the number of fitted
pseudo-samples per fold from
$|\mathcal K_\lambda|\cdot|\Lambda|\,B_{\mathrm{cv}}$ to
$|\Lambda|\,B_{\mathrm{cv}}$ (Supplementary Material, Section~S2).

\subsection{Bias-Corrected Extrapolation}\label{sec:ridge}

Even with optimal $(\lambda_{\max},p_{\mathrm{deg}})$, the polynomial
extrapolation to $\lambda=-1$ may suffer from approximation bias. We
use a two-stage procedure. Stage~1 fits the primary degree-$d$
polynomial $p(\lambda)=\bphi(\lambda)^\top\theta$ with
$\bphi(\lambda)=(1,\lambda,\dots,\lambda^d)^\top$ to the
$(\lambda_j,\bar\beta(\lambda_j))$ points. Stage~2 estimates a
residual correction by ridge regression of the primary-fit residuals
on a richer basis $\bphi_{d+2}(\lambda)$:
\begin{equation}\label{eq:ridge}
  \hat r_{\mathrm{ridge}}(\lambda)=\bphi_{d+2}(\lambda)^\top
  \argmin_{\theta_r}\Bigl\|\bar\beta-\bm P\bar\beta
  -\bPhi_r\theta_r\Bigr\|^2
  +\gamma\|\theta_r\|^2,
\end{equation}
where $\bm P$ is the projection onto the span of the primary basis,
$\bPhi_r$ the $(K\times(d+3))$ basis matrix, and $\gamma=0.01$; the
final estimate is
$\hat\beta_{\asimex}=p(-1)+\hat r_{\mathrm{ridge}}(-1)$. Because both
stages are linear in the curve, the estimate is a single linear
functional $\hat\beta_{\asimex}=\sum_jd_j\bar\beta(\lambda_j)$ with
combined weights $d=c+c_{\mathrm{ridge}}$ satisfying
$\sum_jd_j=1$, the form in which the variance estimator below and
the asymptotic theory analyze it. The ridge stage is $O_p(B^{-1/2})$
at the $\sqrt n$ scale (the theory treats it exactly through the
combined weights).

\subsection{Variance Estimation and Confidence Intervals}\label{sec:var}

Let $c=(c_0,\dots,c_{K-1})$ denote the combined extrapolation weights
(primary polynomial plus ridge second stage), so that the A-SIMEX
estimator is exactly the linear functional
$\hat\beta_{\asimex}=\sum_jc_j\bar\beta(\lambda_j)$ with
$\sum_jc_j=1$, and let
$\bm H_j=\E\bigl[e^{\beta(\lambda_j)^\top W(\lambda_j)}
W(\lambda_j)W(\lambda_j)^\top\bigr]$ be the per-observation Hessian of
the population SIMEX objective at grid point $\lambda_j$. Under the
conditions of Theorem~\ref{thm:an},
\begin{equation}\label{eq:sandwich}
  \Var\bigl(\hat\beta_{\asimex}\bigr)\;\approx\;\frac1n\,\bm C
  \bigl(\btV^{\mathrm{data}}+\tfrac1B\,\btV^{\mathrm{sim}}\bigr)
  \bm C^\top,
  \qquad
  \bm C=\bigl[c_0\bm H_0^{-1},\dots,c_{K-1}\bm H_{K-1}^{-1}\bigr],
\end{equation}
where $\btV^{\mathrm{data}}$ is the covariance across observations of
the stacked conditional influence functions
$\E[\bm m_{ib}(\lambda_j)\mid Y_i,\bw_i]$ and
$\btV^{\mathrm{sim}}$ their conditional covariance given the data;
both admit closed-form expressions under the Gaussian error model
(Gaussian tilt identities) and are estimated from the same replicates
that produce the SIMEX curve, at negligible additional cost.
Two features distinguish \eqref{eq:sandwich} from the additive
``parameter-plus-extrapolation'' split that appears in earlier SIMEX
treatments: (i)~the naive Fisher information
$\bm I_W(\bet_0)^{-1}$ does \emph{not} enter additively: the
estimator is a fixed weighted average of the grid-point
pseudo-estimators, and its variance is the propagated influence
covariance \eqref{eq:sandwich}, which reduces to the naive variance
only in the limit of perfectly correlated grid-point influence; and
(ii)~the simulation noise enters at exactly $O(B^{-1})$, so $B=100$
makes it negligible relative to the data component, as confirmed
empirically in Section~\ref{sec:sims-val}, where the analytic sandwich
recovers $66$--$91\%$ of the parametric-bootstrap variance
(first-order under-statement, improving with $n$). Confidence
intervals are $\hat\beta_{\asimex,j}\pm
z_{1-\alpha/2}\,\widehat{\mathrm{se}}_j$ from
\eqref{eq:sandwich}; under overdispersion the same sandwich is valid
with robust scores, and Section~\ref{sec:sims-nb} verifies this under
negative-binomial outcomes.
